% Rebuilt draft (2026-11-09) of the paper on the paper-protocol branch.
% Build: from the repository root run
%   bash tools/paper_outputs.sh          # tables, figure, numbers -> paper/draft/generated/
%   python tools/annotator_agreement.py  # results/agreement/summary.json (before the line above)
% then compile this file with llncs.cls and splncs04.bst (Springer LNCS kit)
% in this folder, e.g. latexmk -pdf main.tex.
%
% Every number in the text is \res{key}: keys are defined in numbers.txt and
% filled by tools/paper_numbers.py, so no result is typed by hand. A key
% without a result prints as a red ??key. \verify{...} marks a sentence
% whose direction (which model is better) was written from the pilot runs
% and must be checked against the final numbers; \finaltrue hides the notes.
% \tbd{...} marks a fact that is not in any results file.
\documentclass[runningheads]{llncs}

\usepackage[T1]{fontenc}
\usepackage{graphicx}
\usepackage{booktabs}
\usepackage{amsmath}
\usepackage{amssymb}
\usepackage{xcolor}
\usepackage{xr}
\externaldocument[S-]{supplementary}  % compile supplementary.tex first
\usepackage[hidelinks]{hyperref}

\newif\iffinal
% \finaltrue

\makeatletter
\newcommand{\res}[1]{\ifcsname res@#1\endcsname\csname res@#1\endcsname
  \else\textbf{\textcolor{red}{??#1}}\fi}
\makeatother
\iffinal
  \newcommand{\verify}[1]{}
  \newcommand{\tbd}[1]{#1}
\else
  \newcommand{\verify}[1]{\marginpar{\raggedright\tiny\textcolor{red}{Check: #1}}}
  \newcommand{\tbd}[1]{\textcolor{red}{[#1]}}
\fi
\IfFileExists{generated/numbers.tex}{\input{generated/numbers.tex}}{}

% A generated table, or a note where it will go.
\newcommand{\gentable}[1]{\IfFileExists{generated/#1.tex}{\input{generated/#1.tex}}%
  {\begin{table}[t]\centering\fbox{\textcolor{red}{generated/#1.tex: run tools/paper\_outputs.sh}}\end{table}}}

\begin{document}

\title{Learning from Accepted Masks: Distilling a Click-Adapted Interactive Ensemble into SAM~2 for Skin Lesion Annotation}
\titlerunning{Distilling an Adapted Interactive Ensemble into SAM~2}

\author{Athanasios Kallipolitis\inst{1} \and \tbd{co-authors}}
\authorrunning{A. Kallipolitis et al.}
\institute{\tbd{affiliation} \\ \email{\tbd{email}}}

\maketitle

\begin{abstract}
The extended utilization of interactive segmentation in annotation platforms is transforming the production of pixel-level ground truth for medical images, since a few clicks on a promptable model replace the manual delineation of every structure. However, promptable foundation models trained on natural images require several corrective clicks on dermoscopic images, whereas the models that are adapted to the domain either rely on expert masks for their training or are evaluated on the very challenge sets they have been trained on. The paper investigates whether an interactive ensemble of SimpleClick and SAM~2, adapted online on the clicks of a short labelling session, can be distilled into a single SAM~2 model by means of low-rank adapters, using as targets only the masks that a labeller would have accepted from the ensemble without any expert review. The students are compared with general and medical interactive models on the ISIC~2018 test set and on a challenge-disjoint test set drawn from IMA++, with robot and realistic simulated clicks, box prompts and paired bootstrap statistics. The student trained on accepted masks requires fewer clicks than its own teacher and the publicly available models, approaching the student trained on expert masks, which renders the annotation platform capable of improving its own assistant without additional expert effort.\verify{abstract: accepted student better than teacher and all public models on both sets; close to the expert student}
\keywords{Interactive segmentation \and Knowledge distillation \and SAM~2 \and Skin lesion segmentation \and Annotation}
\end{abstract}

\section{Introduction}

The availability of large and accurately annotated image collections constitutes the main prerequisite for the development of supervised deep learning models in medical image analysis, while the pixel-level delineation of lesions and anatomical structures remains one of the most time-consuming and costly steps of the whole process, since it demands the attention of trained experts for every single image. In this context, interactive segmentation has emerged as a practical compromise between fully manual annotation and fully automatic prediction, where a model proposes a mask from a small number of user interactions (positive and negative clicks, boxes or scribbles) and the user corrects it with additional interactions until the result is acceptable \cite{xu2016dios,sofiiuk2022ritm,liu2023simpleclick}. The efficiency of such a scheme is commonly quantified by the number of clicks that is required to reach a target overlap with the reference mask, following the simulated-user protocol introduced in \cite{xu2016dios}, which places every click at the center of the largest remaining error.

Moving a step forward, the introduction of promptable foundation models, namely the Segment Anything Model (SAM) \cite{kirillov2023sam} and its successor SAM~2 \cite{ravi2025sam2}, has provided general-purpose interactive segmenters that are trained on more than a billion masks of natural images and generalize to a wide range of unseen objects. Nevertheless, the experimental studies on medical images have shown that the performance of these models degrades considerably when the appearance of the target differs from that of natural objects, as in the case of low-contrast boundaries and heterogeneous textures \cite{mazurowski2023samstudy}. Towards this direction, a number of medical variants have been proposed: in \cite{ma2024medsam}, the authors fine-tuned SAM on more than a million medical image and mask pairs for box prompts (MedSAM), in \cite{wu2023medsa}, adapters are inserted into SAM in order to specialize it to medical tasks, whereas in \cite{cheng2025imis}, the authors trained an interactive model (IMIS-Net) on a large benchmark of simulated medical interactions. Dermoscopy constitutes a characteristic example of this setting, as the lesions present fuzzy borders, hair, air bubbles, calibration stickers and large variation in size and pigmentation, while the publicly available challenge datasets of the International Skin Imaging Collaboration (ISIC) \cite{codella2019isic,tschandl2018ham} are included, to a large extent, in the training sets of the aforementioned medical models.

Two practical obstacles arise, therefore, when an annotation platform is to be equipped with an interactive assistant for skin lesions. Firstly, the reported performance of medical models on the ISIC challenge test sets is not representative of their performance on new images, since the test images, or images of the same lesions, may have been part of their training data. Secondly, the adaptation of a model to the domain normally requires expert masks, which are precisely the scarce resource that the platform aims to produce. An alternative source of supervision lies in the annotation process itself: every time a labeller accepts a mask proposed by the assistant, a new training pair is created without any additional cost. In this framework, the assistant can be adapted online on the clicks of the current session, and the knowledge of a stronger but slower assistant, such as an ensemble of interactive models, can be transferred to a single and faster model by means of knowledge distillation \cite{hinton2015distill}. The idea of the ensemble derives from the human practice of seeking the opinion of more than one expert before a decision is taken, and its cost at annotation time, where every click has to be processed by every member, is exactly what distillation aims to remove.

In this paper, we propose a distillation scheme in which an interactive ensemble of SimpleClick \cite{liu2023simpleclick} and SAM~2 \cite{ravi2025sam2}, adapted online on a short labelling session, labels a pool of dermoscopic images under a simulated user, and a single SAM~2 student is trained by means of low-rank adapters (LoRA) \cite{hu2022lora} on the prompt encoder and mask decoder to reproduce either the soft maps of the teacher, the masks that the labeller would have accepted from the teacher without expert review, or, as an upper control, the expert masks. Our contribution lies in three aspects. Firstly, we show that the accepted masks, which come for free during annotation, suffice to train a student that requires fewer clicks than its own teacher and approaches the student trained on expert masks.\verify{intro: accepted student beats teacher, close to expert student} Secondly, we introduce an evaluation protocol that accounts for training-set contamination, by complementing the ISIC~2018 test set with a challenge-disjoint test set drawn from the multi-annotator IMA++ collection \cite{abhishek2024imapp}, and for the idealized nature of the robot clicker, by adding realistic simulated clicks from the RClicks clickability model \cite{antonov2024rclicks} and box prompts, with every comparison reported by means of paired bootstrap confidence intervals. Thirdly, we compare the students with general (SAM, SAM~2, SimpleClick, RITM) and medical (MedSAM, IMIS-Net) interactive models, and we relate their accuracy to the agreement between human annotators of the same lesions. The remainder of the paper is organized as follows. In Section~\ref{sec:methods}, the datasets, the models, the distillation scheme and the evaluation protocol are described, whereas Section~\ref{sec:results} presents the experimental results. Section~\ref{sec:discussion} discusses the findings and the limitations of the study, and Section~\ref{sec:conclusions} concludes the paper.

\section{Materials and Methods}\label{sec:methods}

\subsection{Datasets and Test Sets}

The experiments are conducted on dermoscopic images with binary lesion masks, organized into two test sets of 1000 images each and one development set, whose attributes are summarized below. Every image is downscaled so that its longer side does not exceed 1024 pixels, with area interpolation for the images and nearest-neighbor interpolation for the masks.
\begin{itemize}
\item \textbf{ISIC~2018 Task~1} \cite{codella2019isic}: the official training (2594 images), validation (100) and test (1000) splits are used. Twenty training images, drawn with a fixed seed, form the \emph{session} split on which online adaptation takes place, the remaining 2574 training images form the distillation pool, the validation split is used only for checkpoint selection, and the test split only for the final measurements.
\item \textbf{Challenge-disjoint test set}: since IMIS-Net was trained on 90\% of ISIC 2016--2018 Task~1 and PH$^2$ \cite{mendonca2013ph2}, and MedSAM pooled dermoscopic images with a random split, both models have probably encountered the ISIC~2018 test images during training. A second test set is therefore drawn from IMA++ \cite{abhishek2024imapp}, a collection of ISIC Archive images with masks from multiple expert annotators, from which every image that belongs to an ISIC 2016, 2017 or 2018 challenge set (segmentation and classification tasks, the latter including HAM10000 \cite{tschandl2018ham}) is removed, together with every image that shares a lesion or patient identifier with one of them. Out of the \tbd{N} remaining images, 1000 are drawn with a fixed seed, and the reference mask is the STAPLE consensus \cite{warfield2004staple} of the annotators where more than one is available and the single annotator's mask otherwise. Approximately 70\% of these images originate from a single acquisition source, identified by their original resolution of $3024\times2016$ pixels, and their median lesion covers about 2\% of the image.\verify{70\% dominant source and 2\% median lesion size on the final 1000-image set}
\item \textbf{Multi-annotator set}: every challenge-disjoint image with masks from two or more annotators (\res{agr-n} images) is used to estimate the agreement between human annotators.
\end{itemize}

\subsection{Interactive Models and Online Adaptation}

Six interactive models are considered as baselines, all of them in their publicly released form: RITM with an HRNet-18 backbone trained on COCO and LVIS \cite{sofiiuk2022ritm}, the same architecture fine-tuned on the ISIC~2018 training split, SimpleClick with a ViT-Base backbone \cite{liu2023simpleclick}, SAM with a ViT-B encoder \cite{kirillov2023sam}, SAM~2.1 with the Hiera-Tiny encoder \cite{ravi2025sam2}, and IMIS-Net \cite{cheng2025imis}, which is run unmodified from the authors' code. MedSAM \cite{ma2024medsam} is evaluated with box prompts only, since its public checkpoint was trained exclusively on them, following the preprocessing of the authors' inference script.

Online adaptation follows the workflow of the annotation platform, in which the labeller segments a few objects and then requests the adaptation of the assistant on the interactions recorded so far. In this manner, low-rank adapters \cite{hu2022lora} are attached to the parts of each model that are closest to the output (the prompt encoder and mask decoder of SAM and SAM~2, the neck and head of SimpleClick and the OCR head of RITM), while the image encoder remains frozen. Each adapted weight matrix $W$ is replaced by
\begin{equation}
W' = W + \frac{\alpha}{r} B A,
\end{equation}
where $W \in \mathbb{R}^{d\times k}$ is the frozen pretrained weight, $B \in \mathbb{R}^{d\times r}$ and $A \in \mathbb{R}^{r\times k}$ are the trainable low-rank factors, $r$ is the rank (8 for SAM and SAM~2, 4 for SimpleClick and RITM) and $\alpha = 2r$ is a scaling factor. During the session, the simulated user clicks every one of the 20 session images until an IoU of 0.90 or ten clicks are reached, and the adapters are then trained for three consecutive requests on the recorded clicks, each consisting of 30 steps with a learning rate of $2\times10^{-4}$ for SAM, SAM~2 and RITM, and of 10 steps with $10^{-4}$ for SimpleClick, settings that were selected beforehand on natural images. The objective combines the binary cross-entropy and a soft IoU term,
\begin{equation}
\mathcal{L}(z, t) = \mathrm{BCE}\big(\sigma(z), t\big) + 1 - \frac{\sum_i \sigma(z_i)\, t_i}{\sum_i \big(\sigma(z_i) + t_i - \sigma(z_i)\, t_i\big)},
\label{eq:loss}
\end{equation}
where $z$ denotes the predicted logits, $\sigma$ the sigmoid function, $t \in [0,1]$ the target and $i$ runs over the pixels.

\subsection{Teacher, Accepted Masks and Student}

The teacher is an ensemble of SimpleClick and SAM~2, each adapted online on the session split as described above, in which both members receive the same clicks and the output after every click is the per-pixel mean of their probability maps. The choice of these two members is based on preliminary experiments on held-out ISIC~2018 validation images, where they formed the strongest pair among the available models. The teacher labels every image of the distillation pool under the robot clicker: each click is placed at the center of the largest error of the teacher's current mask with respect to the expert mask, until an IoU of 0.95 or five clicks are reached, and the clicks together with the teacher's probability map after each click are stored. It should be noted that the expert masks serve here only as a proxy of the labeller's intent, determining where a click is placed, and never as a training target of the teacher-derived students.

Starting from these records, three kinds of targets are formed for the prediction after the $k$-th click: (i) \emph{teacher}, the soft probability map of the teacher after the same $k$ clicks; (ii) \emph{accepted}, the teacher's final map thresholded at 0.5, i.e., the mask that a labeller would have saved at the end of the interaction without any expert check, used as the target after every click; and (iii) \emph{expert}, the expert mask, which serves as an upper control that the annotation platform would not have at its disposal. The student is a single SAM~2.1 Hiera-Tiny model with LoRA of rank 8 on its prompt encoder and mask decoder, whose image encoder is frozen and whose image features are therefore computed once and cached. At each training step, a record and a click index $k$ are drawn uniformly, the first $k-1$ clicks are replayed without gradients, feeding the low-resolution logits of each prediction to the next one exactly as at annotation time, and the prediction after the $k$-th click is trained towards the corresponding target by means of the loss of Eq.~(\ref{eq:loss}). Training lasts 2000 steps with a batch of four records, AdamW with a learning rate of $2\times10^{-4}$ and a weight decay of $10^{-4}$, a linear warm-up over the first 5\% of the steps followed by cosine decay, and gradient clipping at a norm of 1.0. The adapter is evaluated every 500 steps, including step 0 (the unmodified SAM~2), and the checkpoint with the lowest NoC@90 on the validation split is retained. Each kind of target is trained with three seeds, and the results of a student are averaged over them.

\subsection{Evaluation Protocol}

The main evaluation follows the robot-click protocol \cite{xu2016dios}, where each click is placed at the center (the maximum of the distance transform) of the largest error region, as a positive click on a missed region or a negative click on a false positive, for at most ten clicks per image. The reported measures are the mean IoU after the first click (IoU@1) and the number of clicks needed to reach an IoU of 0.85 and 0.90 (NoC@85, NoC@90), set to ten when the target is never reached. Since the robot clicker is idealized, two further settings are examined. In the first, the clicks are sampled from the clickability map of RClicks \cite{antonov2024rclicks}, a model trained on 475 thousand real clicks on natural images, restricted to the current error region and to either the whole probability mass of the map or its least likely tenth, which emulates an inattentive user. In the second, every model that accepts boxes receives a single box, either tight around the expert mask or with each side shifted by up to 10\% of the box width or height, and IoU and Dice are reported. The random elements of both settings are seeded per image, so that all models receive identical prompts.

The statistical analysis is performed at the level of images. The per-image scores of a student are first averaged over its three seeds, and 95\% confidence intervals of the means are obtained by means of the percentile bootstrap \cite{efron1993bootstrap} with 10{,}000 resamples of the images. Two models are compared by means of the paired bootstrap of their per-image differences, with the two-sided $p$-value defined as twice the smaller share of resampled mean differences lying on either side of zero. Finally, the agreement between annotators is computed on the multi-annotator set as the mean pairwise IoU of the annotators of each image, which provides a reference of the overlap that can reasonably be expected between two valid delineations of the same lesion.

\section{Results}\label{sec:results}

\gentable{click}
\gentable{click_cmp}

The results of the robot-click protocol on both test sets are presented in Table~\ref{tab:click}, while the paired comparisons of NoC@90 are given in Table~\ref{tab:click_cmp}. Starting from the ISIC~2018 test set, the unmodified SAM~2 reaches an IoU@1 of \res{isic-base-iou1} and a NoC@90 of \res{isic-base-noc90}, which improves to \res{isic-adapted-noc90} after online adaptation on the 20 session images. The student trained on the accepted masks reduces NoC@90 to \res{isic-acc-noc90}, a difference of \res{isic-acc-base-d} clicks with respect to SAM~2 (95\% CI \res{isic-acc-base-ci}, \res{isic-acc-base-p}) and of \res{isic-acc-teacher-d} clicks with respect to the student that imitates the soft maps of the teacher (\res{isic-acc-teacher-ci}, \res{isic-acc-teacher-p}).\verify{ISIC: accepted < base, adapted and teacher student in NoC@90} The student trained on the expert masks remains the best of the three, with a NoC@90 of \res{isic-gt-noc90}, leaving a gap of \res{isic-gt-acc-d} clicks (\res{isic-gt-acc-ci}, \res{isic-gt-acc-p}) that represents the price of training without any expert review. Among the publicly available models, IMIS-Net achieves a NoC@90 of \res{isic-imis-noc90} and the RITM model fine-tuned on ISIC~2018 one of \res{isic-ritm-noc90}, both of them having been trained on images of the same challenge, while every baseline before and after online adaptation is reported in Supplementary Table~\ref{S-tab:baselines}.

Moving on to the challenge-disjoint test set, all models require more clicks, which reflects the higher share of small lesions and the different acquisition conditions. The NoC@90 of SAM~2 increases by \res{shift-base-noc90} clicks with respect to ISIC~2018 and that of IMIS-Net by \res{shift-imis-noc90}, whereas the accepted-mask student, with a NoC@90 of \res{clean-acc-noc90}, increases by \res{shift-acc-noc90}.\verify{clean set: IMIS-Net degrades more than the student (contamination argument)} On this set, the accepted-mask student requires \res{clean-acc-base-d} clicks with respect to SAM~2 (\res{clean-acc-base-ci}, \res{clean-acc-base-p}), \res{clean-acc-sc-d} with respect to SimpleClick (\res{clean-acc-sc-p}) and \res{clean-acc-imis-d} with respect to IMIS-Net (\res{clean-acc-imis-p}), while the gap to the expert-mask student amounts to \res{clean-gt-acc-d} clicks (\res{clean-gt-acc-ci}, \res{clean-gt-acc-p}). In this fashion, the advantage of distillation observed on the in-distribution test set is retained, to a certain extent, on images that none of the models has been trained on.\verify{clean set: accepted student better than SAM 2, SimpleClick and IMIS-Net} The evolution of the mean IoU with the number of clicks is illustrated in Fig.~\ref{fig:curves}, where the students reach a higher IoU after the first click, with IoU@1 of \res{clean-acc-iou1} for the accepted-mask student against \res{clean-base-iou1} for SAM~2 on the challenge-disjoint set, and the curves converge after a few corrective clicks.\verify{curves: students higher at click 1, converging later}

\begin{figure}[t]
\centering
\IfFileExists{generated/curves.pdf}{\includegraphics[width=\textwidth]{generated/curves.pdf}}%
  {\fbox{\textcolor{red}{generated/curves.pdf: run tools/paper\_outputs.sh}}}
\caption{Mean IoU against the number of robot clicks on the ISIC~2018 and the challenge-disjoint test sets. Student curves average three seeds.}
\label{fig:curves}
\end{figure}

\gentable{robust_clean}

The robustness of the models to the type and the quality of the prompt on the challenge-disjoint test set is presented in Table~\ref{tab:robust_clean}, while the corresponding results on ISIC~2018 are given in Supplementary Table~\ref{S-tab:robust_isic} and the separate box and realistic-click tables, with Dice and IoU@1, in Supplementary Tables~\ref{S-tab:box_isic}--\ref{S-tab:rclicks_clean}. With a tight box, the accepted-mask student reaches an IoU of \res{clean-box-acc-iou}, against \res{clean-box-sam2-iou} for SAM~2, \res{clean-box-medsam-iou} for MedSAM and \res{clean-box-imis-iou} for IMIS-Net, with differences of \res{clean-box-acc-medsam-d} (\res{clean-box-acc-medsam-p}) and \res{clean-box-acc-imis-d} (\res{clean-box-acc-imis-p}) with respect to the two medical models, although the student has been trained on clicks only.\verify{box: student on par with or above MedSAM and IMIS-Net on the clean set} When the sides of the box are shifted by up to 10\%, the IoU of the student changes by \res{clean-jit-acc}, that of SAM~2 by \res{clean-jit-sam2}, that of MedSAM by \res{clean-jit-medsam} and that of IMIS-Net by \res{clean-jit-imis}. Returning to clicks, the realistic clicks of RClicks increase the NoC@90 of the accepted-mask student to \res{clean-rc-acc-noc90}, i.e., by \res{clean-rc-acc-robot-d} clicks with respect to the robot clicker (\res{clean-rc-acc-robot-p}), whereas the difference from SAM~2 under the same clicks is \res{clean-rc-acc-sam2-d} (\res{clean-rc-acc-sam2-p}). When the clicks are drawn from the least likely tenth of the clickability map, the NoC@90 of the student and of SAM~2 rise to \res{clean-rclo-acc-noc90} and \res{clean-rclo-sam2-noc90}, respectively, with a difference of \res{clean-rclo-acc-sam2-d} (\res{clean-rclo-acc-sam2-p}), which indicates that the benefit of the student is not an artifact of the idealized placement of the robot clicks.\verify{RClicks: student still better than SAM 2 under realistic and least likely clicks}

Moreover, the challenge-disjoint test set is examined by subgroups, as shown in Supplementary Table~\ref{S-tab:subgroups}. On the images of the dominant acquisition source, the NoC@90 of the accepted-mask student is \res{dom-acc-noc90} against \res{dom-base-noc90} for SAM~2, while on the remaining sources the respective values are \res{oth-acc-noc90} and \res{oth-base-noc90} (difference \res{oth-acc-base-d}, \res{oth-acc-base-p}), and on lesions covering at least 5\% of the image \res{ge5-acc-noc90} and \res{ge5-base-noc90}.\verify{subgroups: the student's advantage holds outside the dominant source} Finally, the agreement between human annotators on the \res{agr-n} multi-annotator images provides a reference for these results: the mean pairwise IoU is \res{agr-mean} (95\% CI \res{agr-lo}--\res{agr-hi}, median \res{agr-median}), and a second annotator reaches an IoU of 0.85 with the first on only \res{agr-ge85} of the images and of 0.90 on \res{agr-ge90}. Consequently, a target of 0.90 lies, for a considerable share of the lesions, beyond the overlap between two experts, and the NoC@90 differences between the strongest models should be interpreted in this light.\verify{agreement: second annotator below 0.90 on a large share of images}

\section{Discussion}\label{sec:discussion}

The main goal of the study was to examine whether an annotation platform can improve its own interactive assistant from the masks that its users accept, without resorting to additional expert annotation. In this context, the results indicate that a single SAM~2 model, distilled from an online-adapted ensemble of SimpleClick and SAM~2, requires fewer clicks than its teacher and than the general and medical interactive models that were examined, both on the ISIC~2018 test set and on images outside every ISIC challenge.\verify{discussion: summary of the main finding} The fact that the student trained on accepted masks outperforms the student that imitates the soft maps of its teacher can be attributed to the nature of the target: the accepted mask summarizes the end result of a whole interaction, so that already after the first click the student is pushed towards the final mask instead of the teacher's intermediate and uncertain predictions. Furthermore, the student processes every click with a single model, which reduces the latency of the assistant with respect to the two-member teacher, an important property for the responsiveness of an annotation tool.

An important limitation of the study is that the users are simulated. The robot clicker derives the position of each click from the expert mask, which implies that the accepted masks, although never compared with an expert during training, are produced by clicks that are guided by expert knowledge; in practice, a labeller would stop at a mask that is visually acceptable rather than at a fixed IoU. To a certain extent, this is mitigated by the realistic clicks of RClicks and by the box prompts, under which the ranking of the models is preserved,\verify{discussion: ranking preserved under RClicks and boxes} although the clickability model was trained on natural images and does not describe the behavior of dermatologists. Moreover, the exclusion of challenge images from the second test set relies on the identifiers of the ISIC Archive, and since most images carry no lesion or patient identifier, the check for lesions shared with the challenge sets is only partial. The second test set is also dominated by a single acquisition source and by small lesions, which is why the results are reported by subgroup as well. Finally, the reference masks themselves are uncertain, as the moderate agreement between annotators shows, and a user study with dermatologists on the annotation platform would be required in order to translate click savings into annotation time.

\section{Conclusions}\label{sec:conclusions}

In this work, we have investigated the distillation of an online-adapted interactive ensemble into a single SAM~2 model for the annotation of skin lesions, using as targets the masks that a labeller would accept from the ensemble without expert review. The evaluation on the ISIC~2018 test set and on a challenge-disjoint test set, with robot and realistic clicks, box prompts and paired bootstrap statistics, showed that such a student is on par with or better than the publicly available general and medical interactive models and approaches a student trained on expert masks, while the agreement between human annotators places the remaining differences into perspective.\verify{conclusions: summary of the findings} Therefore, an annotation platform can use its own accepted masks to improve its assistant as annotation proceeds. Future work should be directed towards the validation of the scheme with real labellers and annotation times, the repetition of the distillation over successive annotation rounds, and its extension to other imaging modalities and annotation tasks of the platform.

\bibliographystyle{splncs04}
\bibliography{references}

\end{document}
